\newcommand{\conone}{\\
The Java program successfully sorts the first half of the array in ascending order and the second half in descending order based on user input. It demonstrates the effective use of array manipulation, Java’s built-in \texttt{Arrays.sort()} method, and \\\texttt{Collections.reverseOrder()}. This task strengthens the understanding of partial array sorting and highlights the importance of handling primitive and wrapper types in Java. It also reinforces logical thinking and modular programming practices.
}

\newcommand{\contwo}{\\
The Java program efficiently reads a matrix from user input and computes the minimum values for each row and each column. It demonstrates the use of nested loops and conditional logic to traverse two-dimensional arrays. This assignment enhances understanding of matrix manipulation and reinforces key concepts such as array indexing, comparison operations, and formatted output in Java. It serves as a practical example of solving numerical problems using structured programming techniques.
}

\newcommand{\conthree}{\\
The Java program effectively removes all consonants from the input string while preserving vowels and non-letter characters. It utilizes character-level iteration, conditional checks, and string manipulation through the \texttt{StringBuilder} class. This assignment reinforces the understanding of string handling, method creation, and character classification in Java. It also highlights the importance of clean logic and efficient looping in solving text processing problems.
}

\newcommand{\confour}{\\
The Java program successfully demonstrates the implementation of a class to represent a 2D point, including constructors, accessor and mutator methods, and a method to calculate the distance between points. Through the \texttt{MyPoint} and \texttt{TestMyPoint} classes, the assignment reinforces core object-oriented concepts such as encapsulation, method overloading, and the use of instance methods. It provides hands-on experience in class design and working with coordinate geometry in Java.
}

\newcommand{\confive}{\\
The Java program demonstrates the concept of inheritance by creating a superclass \texttt{Person} and extending it into subclasses \texttt{Student} and \texttt{Staff}. It effectively showcases constructor chaining using \texttt{super}, method overriding through the \texttt{toString()} method, and subclass-specific properties and methods. This assignment strengthens the understanding of object-oriented principles, particularly inheritance and polymorphism, enabling the design of reusable and organized class hierarchies in Java.
}

\newcommand{\consix}{\\
The Java program demonstrates the use of inheritance and method overriding by modeling geometric shapes through the \texttt{Square} and \texttt{Cylinder} classes. The \texttt{Square} class is reused to represent a cube, while the \texttt{Cylinder} class extends it to calculate cylindrical volume using overridden behavior. This assignment highlights the importance of polymorphism, constructor chaining with \texttt{super}, and runtime method resolution. It reinforces how base classes can be extended and customized to solve more specific problems in object-oriented programming.
}

\newcommand{\conseven}{\\
The Java program effectively demonstrates the concept of interfaces through the \texttt{Engine} interface and its implementation in the \texttt{VehicleEngine} class. It showcases how abstraction is achieved in Java by defining behavior contracts that must be implemented by concrete classes. The program also illustrates real-world functionality like speeding up and changing gears, helping to understand how interface-based design supports modular, flexible, and scalable software development.
}

\newcommand{\coneight}{\\
The Java program clearly demonstrates the concept of exception handling using \texttt{try-catch} blocks. It handles both \\\texttt{ArrayIndexOutOfBoundsException} and \texttt{ArithmeticException}, showing how runtime errors can be gracefully managed without abruptly terminating the program. This assignment enhances understanding of robust programming practices, emphasizing the importance of error detection, exception handling, and ensuring program continuity in Java applications.
}

\newcommand{\connine}{\\
The Java program successfully demonstrates the use of JDBC (Java Database Connectivity) to perform basic database operations such as inserting and retrieving student records from a MySQL database. It covers essential database interactions using \texttt{Connection}, \texttt{PreparedStatement}, \texttt{Statement}, and \texttt{ResultSet} classes. This assignment enhances understanding of database-driven application development and reinforces concepts of secure data handling, exception management, and SQL integration within Java programs.
}

\newcommand{\conten}{\\
The Java program effectively demonstrates how to remove a specified element from an array by creating a new array and omitting the target element. It utilizes conditional logic and loop control to handle array traversal and element shifting. This assignment reinforces the understanding of array manipulation, user input handling, and memory management through the use of dynamic array reconstruction. It also highlights the importance of validating user inputs and providing appropriate feedback.
}

\newcommand{\coneleven}{\\
The Java program successfully demonstrates how to insert an element at a specific position in an array. It includes input validation, element shifting, and efficient array traversal. This assignment helps in understanding core concepts of array manipulation, boundary checking, and memory management in Java. It reinforces logical thinking, precise index handling, and user-driven data insertion techniques in a structured programming approach.
}

\newcommand{\contwelve}{\\
The Java program efficiently identifies all pairs of elements in an array whose sum matches a user-defined target value. It demonstrates the use of nested loops, conditional checks, and boolean flags to handle pairwise comparison. This assignment strengthens the understanding of array traversal, combination logic, and user input handling. 
}

\newcommand{\conthirteen}{\\
The Java program effectively demonstrates how to remove duplicate elements from an array using the \texttt{LinkedHashSet} class, which maintains insertion order while eliminating repeated values. This approach simplifies the logic for handling duplicates and ensures clean, concise output. The assignment enhances the understanding of Java Collections, set operations, and array-to-collection conversions. It also reinforces concepts of efficient data handling and memory management in real-world programming tasks.
}

\newcommand{\confourteen}{\\
The Java program efficiently calculates the length of the longest consecutive sequence within an unsorted array using a \texttt{HashSet} for constant-time lookups. By checking only for the start of sequences and extending them conditionally, the solution avoids unnecessary comparisons and improves performance. This assignment strengthens the understanding of set-based algorithms, loop optimization, and problem-solving using hashing techniques in Java.
}

\newcommand{\confifteen}{\
This Java program demonstrates lexicographical comparison of two strings using the \texttt{compareTo()} method. It takes two strings as input and compares them based on Unicode values of characters. The result determines whether the first string is equal to, less than, or greater than the second string. This reinforces concepts of string manipulation, conditional branching, and understanding lexicographic order in Java.
}

\newcommand{\consixteen}{\\
The Java program demonstrates the use of the \texttt{regionMatches()} method to compare specific regions (substrings) of two strings. It provides insight into how substring-level comparisons can be efficiently performed without extracting substrings manually. This assignment enhances the understanding of Java's built-in string manipulation methods, and their role in implementing precise and performance-optimized text comparison operations.
}

\newcommand{\conseventeen}{\\
The Java program successfully generates all possible permutations with repetition for a given string using recursion. It demonstrates the use of backtracking and string concatenation within a recursive method to build each permutation. This assignment enhances the understanding of recursive problem-solving techniques, especially in combinatorics, and provides a foundation for exploring more complex algorithms involving string manipulation and enumeration.
}

\newcommand{\coneighteen}{\\
The Java program accurately counts the number of words in a given string by using string trimming and regular expression-based splitting. It handles various spacing scenarios, including multiple spaces and leading/trailing whitespace. This assignment strengthens string manipulation skills and introduces basic text processing techniques in Java, which are essential in many real-world applications such as form validation, parsing, and natural language processing.
}

\newcommand{\connineteen}{\\
The Java program illustrates the concept of a copy constructor by initializing one object using another of the same class. It reinforces object-oriented programming principles such as encapsulation and constructor overloading. By copying the properties of an existing Box object, this program demonstrates how to create deep copies of objects, which is a crucial technique in avoiding unintentional reference sharing and maintaining data integrity in real-world applications.
}

\newcommand{\contwenty}{\\
This Java program effectively demonstrates constructor overloading, including the use of a copy constructor, to manage employee salary-related data. It calculates Provident Fund (PF), Allowance, and Total Salary based on a given basic salary and encapsulates this logic within a method. By separating initialization, computation, and display, the program maintains clean structure and readability, reflecting good object-oriented design practices.
}

\newcommand{\contwentyone}{\
This Java program demonstrates the use of custom exception handling by creating a user-defined exception class \texttt{NegativeSizeException}. It checks whether an array size is negative before creation and throws the custom exception if so. This approach shows how to enhance error reporting and program robustness by handling specific application-level errors more meaningfully.
}

\newcommand{\contwentytwo}{\
The Java program demonstrates object comparison by defining a static method \texttt{isEqual()} in the \texttt{Student} class. It compares two \texttt{Student} objects based on both roll number and name, illustrating how logical equality can be implemented beyond reference equality. This exercise reinforces the use of parameterized constructors, static utility methods, and fundamental principles of encapsulation and data comparison in object-oriented programming.
}

\newcommand{\contwentythree}{\
This Java program illustrates the use of abstract classes and method overriding in an employee management scenario. The abstract class \texttt{Employee} defines a base structure with a constructor, an abstract method \texttt{calculateNetSalary()}, and a concrete \texttt{display()} method. Subclasses \texttt{Manager} and \texttt{Clerk} override the abstract method to provide their own net salary calculations using different allowance and deduction percentages. This demonstrates runtime polymorphism, abstraction, and clean class hierarchy design.
}